\paragraph{\pal.} A write-time belief store. Every observed key holds an append-only list of versions (belief time, admitted report ids, interpretation set computed with the declared semantics); every derived belief (entity, derived attribute) holds versions that \emph{pin} the key versions it was justified by. Answers at belief time $\tau$ are read from the versions in force at $\tau$, so bitemporal questions are lookups. On an assertion, the key is recomputed and every dependent derived belief is refreshed (ordinary belief updating). On a retraction or withdrawal, the key is recomputed and---this is the propagation step---every derived belief that depends on it is refreshed. Support (\S\ref{sec:support}) is computed from the admitted evidence of a key. In the natural-language setting the store reasons over the shared extracted claims; the LLM is used only at the boundary (extraction), never to answer.

\paragraph{Ablations.} \emph{No propagation}: a retraction updates the retracted key only; derived beliefs keep their pinned versions until an ordinary assertion refreshes them. \emph{Forced commit}: each key keeps a single interpretation, the admissible one with fewest error labels (ties: ended over continuing). Further ablations of the support model are reported with RQ3.

\paragraph{Symbolic replay (the equivalence baseline).} The same claims in an append-only log; at query time the semantics is executed from scratch over the log filtered to the query's belief time: every key the query depends on is re-enumerated and the derived value re-chained, with the per-key cache cleared before each query (\emph{cold}) or retained across queries (\emph{cached}): the cache is a memo keyed by (key, belief time, policy) that is never invalidated, which is sound because the log is append-only and a report at a later belief time cannot change an earlier belief time's interpretation set; it therefore helps only when a later query touches the same key at the same belief time. Replay and \pal{} execute the same declared semantics on the same evidence and have no information the other lacks; they must agree on every answer, and \S\ref{sec:res-systems} verifies that they do and measures what each costs. In the natural-language settings replay is the ``log + rule adjudicator'' row: what query-time reasoning achieves over the extracted claims when the reasoning is correct.

\paragraph{Append-only log with LLM adjudication.} The same extracted claims in the same log; at query time the backbone LLM receives the log filtered to the query's belief time, the declared attribute and rule block, the shared instruction, and the query rendered as its answer slot, and must return the contract JSON, with one repair attempt. Context-budget variants restrict the log to the $k$ most recent entries plus every entry on the queried key ($k\in\{10,20,40\}$); a raw-report variant adjudicates the natural-language reports directly, with no extraction. Relative to \pal{} this comparison changes two things at once, who reasons (a prompted model instead of a rule engine) and when (query time instead of write time); replay, which reasons with the rule engine at query time, separates them.

\paragraph{Temporal DB (last write wins).} A rule-based bitemporal store without justifications: the record with the latest valid-time anchor wins, retractions remove records, no error hypotheses, derived values computed at query time. It differs from \pal{} in two respects at once, no justifications and a commit-always uncertainty policy, so its gap is not attributed to one of them; the always-commit reference (latest report by report time) isolates pure recency.

\paragraph{Support model.}\label{sec:support} Support is the probability that a candidate value is true for its interval, defined from admitted evidence only: evidence is grouped by origin (copies count once), each origin contributes a per-class log-odds (declared, or learned from verified outcomes of training streams), discounted by staleness on changeable attributes with a 20-day half-life, boosted by a change cue, and capped for a report targeted by a correction at half the correcting origin's weight (the min(provenance, content) cap); support of $v$ is the logistic of the net evidence for $v$ against all competitors, with a prior offset when no competitor exists. Source identities are never used, only classes and origins. Support is produced for exactly the alternatives the justified-belief layer keeps open, so the two layers are evaluated separately.
